\documentclass[10pt,a4paper]{amsart}
\usepackage[margin=19mm]{geometry}
\usepackage{amsmath,amssymb,mathtools}
\usepackage[scr=rsfs,cal=euler]{mathalfa}
\usepackage{xfrac}
\usepackage[unicode,hidelinks,bookmarks=false]{hyperref}
\newcommand{\ie}{i.e.\ }
\newcommand{\cf}{cf.\ }
\newcommand{\resp}{respectively\ }
\newcommand{\Subsection}{\subsection}
\providecommand{\nobreakdash}{\nobreak\hbox{-}}
\setlength{\emergencystretch}{2em}
\multlinegap0pt
\newtheorem{lemma}{Lemma}
\title{Curved Kakeya problems and \\the projective geometry of paths}
\author{Shaoming Guo, Larry Guth, Arian Nadjimzadah, Minxing Shen, Ruixiang Zhang}
\date{Source: arXiv:2607.27629v1 (CC BY 4.0)}
\begin{document}
\maketitle

\noindent\emph{Source and licence.} Curved Kakeya problems and \\the projective geometry of paths. Version arXiv:2607.27629v1 (CC BY 4.0), distributed by arXiv under the Creative Commons Attribution 4.0 International licence, http://creativecommons.org/licenses/by/4.0/, as recorded in the arXiv licence metadata for this exact version and shown on the abstract page https://arxiv.org/abs/2607.27629v1. Full licence notice: \texttt{licenses/arxiv-2607.27629-CC-BY-4.0.txt}.\par\smallskip

\noindent\textit{Editorial scope.} Counted unit: the article's lemma lemma (source0.tex lines 3325--3334). The article's complete original proof of it is delivered with it (source0.tex lines 3336--3422, one \texttt{proof} environment of 87 lines). No mathematics is written, repaired or re-derived: the statement and the proof are the article's own lines, in the article's own notation, and no retained line is substituted or re-flowed.  No further source-local context is needed: every cross-reference of every retained line resolves inside the two delivered spans.  Nothing was omitted from any retained span, so the package carries no omission record.  No retained line contains a citation, so the wrapper carries no bibliography and no bibliography entry.  No retained line contains a figure environment, an image-inclusion command or drawing code, so no image is delivered and the package declares no asset dependency.  Editorial formatting only, in addition: an AMS \texttt{amsart} preamble that restates the article's own declarations and macros used by the retained lines, namely the environments \texttt{lemma} on separate counters, together with the standard packages those lines need.  The retained environments print in this document as Lemma 1 in order of appearance, whereas the article prints the same items with the numbering of its own full document.  The complete native source of the article as distributed in the arXiv e-print for this exact version is bundled unchanged as \texttt{sources/206361/source0.tex}.
    \begin{lemma}[Linear algebra lemma]\label{lem: lin alg lemma}
    Let $T$ be a $(1,3)$-tensor on 3-dimensional vector space $V$ satisfying the following properties: 
    \begin{enumerate}
        \item (Skew symmetry) $T(X,Y,Z) = -T(Y,X,Z)$ 
        \item (Bianchi identity) $T(X,Y,Z) + T(Y,Z,X)+T(Z,X,Y) = 0$ 
        \item (Trace-free) $\mathrm{trace}(X \mapsto T(X,Y,Z)) = 0$.
        \item (Plane property) $T(X,Y,Y) \in \mathrm{Span}(X,Y)$ for all $X,Y \in V$.
    \end{enumerate}
    Then $T \equiv 0$. 
    \end{lemma}

    \begin{proof}[Proof of Lemma \ref{lem: lin alg lemma}]

    For a basis $X,Y,Z \in V$, write $X^*, Y^*, Z^* \in V^*$ for the dual basis. \\
    
    \noindent \textbf{Step 1: $T(X,Y,Y) \| Y$ and $T(Y,X,X) \| X$}\\
    
    The second claim follows from the first by swapping $X$ and $Y$. We focus on the first claim. 
    By the plane property, write 
    \begin{equation}
    T(X,Y,Y) = aX + bY.
    \end{equation}
    The goal is to show $a = 0$. By the trace-free condition, 
    \begin{align} \label{eq: 1 alg lemma}
        0 &= \mathrm{trace}(U \mapsto T(U,Y,Y)) \\
        &=X^*T(X,Y,Y) + Y^*T(Y,Y,Y) + Z^*T(Z,Y,Y) \\
        &= a + 0 + Z^*T(Z,Y,Y),
    \end{align}
    where we used skew-symmetry for the second summand. Apply the plane property to the independent vectors $X + \epsilon Z, Y$ to get 
    \begin{align}\label{eq: 2 alg lemma}
        T(X+ \epsilon Z,Y,Y) = \lambda(\epsilon) (X + \epsilon Z) + \mu(\epsilon) Y. 
    \end{align}
    Taking $\epsilon = 0$ and applying $X^*$ in \eqref{eq: 2 alg lemma} gives 
    \begin{align}
        a = \lambda(0). 
    \end{align}
    Taking a derivative in $\epsilon$ at $0$ and applying $Z^*$ in \eqref{eq: 2 alg lemma} gives 
    \begin{align}
        Z^*T(Z,Y,Y) = \lambda(0) = a.
    \end{align}
    Plugging this into \eqref{eq: 1 alg lemma}, we find $a = 0$.\\


\noindent \textbf{Step 2: $T(Y,X,X), T(X,Y,Y) = 0$}\\

Write $T(Y,X,X) = aX$ and $T(X,Y,Y) = bY$. We will show $a,b=0$. Define the coefficients 
\begin{align}
    r &= Z^* T(X,Y,Z) \\
    s &= Z^*T(Y,Z,X) \\
    t &= Z^*T(Z,X,Y).
\end{align}
By the plane-property, 
\begin{align}\label{eq: 3 alg lemma}
    T(X, Y+\epsilon Z, Y +\epsilon Z) = \lambda(\epsilon) (Y + \epsilon Z) + \mu(\epsilon) X,
    \end{align}
    which implies $r-t=b$. 
    Applying the plane property again (of course for a possibly different choice of $\lambda, \mu$), 
    \begin{align}
        T(Y, X + \epsilon Z, X + \epsilon Z) = \lambda(\epsilon) (X+ \epsilon Z) + \mu(\epsilon)Y.
    \end{align}
    This implies $s-r = a$. The Bianchi identity gives $r + s + t =0$. We will also need the following two applications of the trace-free condition: 
    \begin{align}
        0 &= \mathrm{trace}(U \mapsto T(U,X,Y)) \\
        &= -b + t \\
        0 &= \mathrm{trace}(U \mapsto T(U,Y,X)) \\
        &= -a - s.
    \end{align}
    In total, all the equations are: 
    \begin{align}
        r - t &= b \\
        s-r &= a \\
        r + s + t &= 0 \\
        b &= t \\
        a &= -s.
    \end{align}
    These equations easily imply $a = b = 0$. \\

    \noindent \textbf{Step 3: $T \equiv 0$}\\

    We now use $T(X,Y,Y) = 0$ for all $X,Y \in V$ to conclude that $T \equiv 0$. By using
    \begin{equation}
    T(X, Y+ Z, Y+Z) =0,
    \end{equation}
     expanding, and using 
     \begin{equation}
     T(X,Y,Y) = T(X,Z,Z) = 0,
     \end{equation}
      we get  
    \begin{align}
        T(X,Y,Z) &= - T(X,Z,Y). 
    \end{align}
    Thus $T$ is skew-symmetric in the second two entries. Combined with the skew-symmetry in the first two entries and the Bianchi identity,
    \begin{align}
        0 &= T(X,Y,Z) + T(Y,Z,X) + T(Z,X,Y) \\
        &= 3T(X,Y,Z),
    \end{align}
    and hence $T \equiv 0$. 
\end{proof}

\end{document}
